\section{Non-Cohen--Macaulay sets}\label{sec:noncm}

This section proves Theorem~\ref{thm:D}. A set that is not CM has fewer sums than its intersection form at finitely many levels. We show that in each residue class these levels form an interval, and that once $hd\ge c_1+c_2$ the total of $I_h$ is exactly the mass, so that every sum missing from $hA_0$ raises the total above it.

For a set $E$ of integers containing $0$ and an integer $n\ge0$ put $D_E(n)=n-\max(E\cap[0,n])$, so that $D_h=D_{hA_0}$. For every $A_0$ and $h\ge0$ we have $hA_0\subseteq I_h$ and $0,hd\in hA_0$. Put $F_h=I_h\setminus hA_0$, and call $h$ a \emph{failing level} if $F_h\neq\varnothing$. By Theorem~\ref{thm:cm}, $A_0$ is CM if and only if it has no failing level.

\begin{proposition}[Failing levels]\label{prop:failing-interval}
Fix $1\le s\le d-1$. The levels $h$ at which some integer $kd-s$ lies in $F_h$ form an interval of integers, possibly empty. Consequently $A_0$ has finitely many failing levels.
\end{proposition}

\begin{proof}
\emph{The corners.} Write a target of class $-s$ as $kd-s$ and use coordinates $(k,\ell)$ with $\ell=h-k$. Let $\mathcal C_s$ be the set of $(k,\ell)\in\N^2$ with $kd-s\in(k+\ell)A_0$, and let $\mathcal C^I_s$ be the set with $kd-s\in I_{k+\ell}$. By Lemma~\ref{lem:residue} with $m=d$, $\mathcal C^I_s=[u_s,\infty)\times[v_s,\infty)$. The set $\mathcal C_s\subseteq\mathcal C^I_s$ is closed under adding $(1,0)$ (add a part $d$) and $(0,1)$ (add a part $0$), so it is the union of the quadrants above its finitely many minimal points (Dickson's lemma) $(\kappa_1,\mu_1),\dots,(\kappa_t,\mu_t)$, ordered so that $\kappa_1<\dots<\kappa_t$ and $\mu_1>\dots>\mu_t$. We claim $\kappa_1=u_s$ and $\mu_t=v_s$. The Ap\'ery element $u_sd-s$ of $S_1$ is a sum of some number $j$ of positive parts, so $u_sd-s\in jA_0$; as all parts are at most $d$ and $u_sd-s>(u_s-1)d$, we have $j\ge u_s$, so $(u_s,j-u_s)\in\mathcal C_s$. Likewise the Ap\'ery element $v_sd+s$ of $S_2$ is a sum of $j'$ parts $d-a$, $a\in A_0$, so $j'd-(v_sd+s)=(j'-v_s)d-s$ lies in $j'A_0$, and $(j'-v_s,v_s)\in\mathcal C_s$. Together with $\mathcal C_s\subseteq\mathcal C^I_s$ this proves the claim.

\emph{The failing levels of one class.} If $t=1$, then $\mathcal C_s=\mathcal C^I_s$ and no level fails in this class. If $t\ge2$, then $\mathcal C^I_s\setminus\mathcal C_s$ is the disjoint union of the rectangles $[\kappa_j,\kappa_{j+1}-1]\times[v_s,\mu_j-1]$, $1\le j<t$, because for $\kappa_j\le k<\kappa_{j+1}$ the least $\ell$ with $(k,\ell)\in\mathcal C_s$ is $\mu_j$. The levels $k+\ell$ met by rectangle $j$ form the interval $J_j=[\kappa_j+v_s,\ \kappa_{j+1}+\mu_j-2]$, which is nonempty. For $j\le t-2$, strict descent gives $\mu_j\ge\mu_{j+1}+1\ge v_s+2$, so the upper end of $J_j$ is at least $\kappa_{j+1}+v_s$, the lower end of $J_{j+1}$. Hence $J_1\cup\dots\cup J_{t-1}$ is the interval $[u_s+v_s,\ \max_{j<t}(\kappa_{j+1}+\mu_j-2)]$. Multiples of $d$ never fail, so the failing levels are the finite union of these intervals over $s$.
\end{proof}

In other coordinates, Proposition~\ref{prop:failing-interval} and the formula for $\tau$ below are \cite[Lemma~2.5]{CompanionRegularity}; we keep the short proof.

Let $\tau$ be the least $h_0\ge0$ such that $hA_0=I_h$ for every $h\ge h_0$. Thus $\tau=0$ for CM sets, and otherwise $\tau-1$ is the last failing level; by the proof above, $\tau-1=\max_s\max_{j<t_s}(\kappa_{j+1}+\mu_j-2)$ over the classes $s$ with $t_s\ge2$ corners, an explicit finite computation.

\begin{proposition}[Overshoot]\label{prop:overshoot}
Let $h\ge0$ satisfy $hd\ge c_1+c_2$. Then $\sum_{n=0}^{hd}D_{I_h}(n)=\M(S_1)+\M(S_2)$. Consequently $\Tot(h)\ge\M(S_1)+\M(S_2)$, with equality if and only if $h$ is not a failing level. In particular, if $A_0$ is not CM and $(\tau-1)d\ge c_1+c_2$, then $\Tot(\tau-1)>\M(S_1)+\M(S_2)$.
\end{proposition}

\begin{proof}
Since $hd-c_1\ge c_2$, the integer $c_1$ lies in $I_h$, and so does every integer of $[c_1,hd-c_2]$. An integer $n<c_1$ has $hd-n>c_2$, so $n\in I_h$ if and only if $n\in S_1$; an integer $n>hd-c_2$ has $n>c_1$, so $n\in I_h$ if and only if $hd-n\in S_2$. Thus the integers of $[0,hd]$ outside $I_h$ are the gaps of $S_1$, all below $c_1$, and the reflections $hd-x$ of the gaps $x$ of $S_2$, all above $hd-c_2$, and the element $c_1$ of $I_h$ separates the two kinds. A maximal run of $L$ gaps of $S_1$ contributes $1+\dots+L$ to the total of $I_h$ and to $\M(S_1)$. For a reflected run, the undershoot at $hd-x$ is the distance from $x$ up to the next element of $S_2$, so by Lemma~\ref{lem:elementary}(iv) the run contributes $L(L+1)/2$, as in $\M(S_2)$. This proves the first identity (if a conductor is $0$, the corresponding set of gaps is empty). Removing from a set containing $0$ and $hd$ an element $x$ with $0<x<hd$ makes its undershoot at $x$ positive and decreases no other undershoot (Lemma~\ref{lem:elementary}(iii)); since $hA_0=I_h\setminus F_h$ and $0,hd\notin F_h$, the total of $hA_0$ exceeds that of $I_h$ exactly when $F_h\neq\varnothing$.
\end{proof}

So once the range is long enough, the intersection form carries exactly the mass, and each sum that $hA_0$ lacks adds to the total. For a CM set Proposition~\ref{prop:overshoot} recovers $\Tot(h)=\M(S_1)+\M(S_2)$ whenever $hd\ge c_1+c_2$. For four-element sets the hypothesis on $\tau$ always holds: by \cite[Theorem~A]{CompanionRegularity}, every non-CM four-element set satisfies
\[
(\tau-1)d\ge c_1+c_2+d ,
\]
with equality exactly for $\{0,b-1,b+1,2b\}$ with $b\ge2$ even.

\begin{corollary}\label{cor:four-overshoot}
Every four-element set $A_0$ that is not CM satisfies $\Tot(\tau-1)>\M(S_1)+\M(S_2)$. In particular, for four-element sets, $\Tot(h)\le\M(S_1)+\M(S_2)$ for all $h$ implies that $A_0$ is CM.
\end{corollary}

\begin{proof}
Combine \cite[Theorem~A]{CompanionRegularity} with Proposition~\ref{prop:overshoot}.
\end{proof}

With five elements the conclusion can fail.

\begin{example}\label{ex:five-element}
Let $A_0=\{0,5,9,11,20\}$, the non-Cohen--Macaulay set of \cite[Example~4.3]{BermejoGimenez2006} (see also \cite[Example~2.13(2)]{GimenezGonzalezSanchez2023}). Then $c_1=18$, $c_2=44$ and $\M(S_1)+\M(S_2)=74$. The only failing level is $h=3$, so $A_0$ is not CM and $\tau=4$, and $\Tot(h)=53,69,74$ for $h=1,2,3$. For $h\ge4$ we have $hA_0=I_h$ and $20h\ge80>c_1+c_2=62$, so $\Tot(h)=74$ by Proposition~\ref{prop:overshoot}. Hence $\Tot(h)\le\M(S_1)+\M(S_2)$ for every $h$: the set never overshoots its mass. At the failing level $h=3$ the hypothesis $hd\ge c_1+c_2$ fails ($60<62$), and the total equals the mass although $F_3=\{24\}$: the total of $I_3$ is $73$, and the correction of Proposition~\ref{prop:correction} at $z=24$, with $z^-=23$ and $z^+=25$, adds $1$.
\end{example}
